\section{Ba phép đo và ba cái nền}
\label{sec:ch11-ba-phep-do-va-ba-cai-nen}

\index{đánh giá!nền của phép đo|(}%
Một tiên đề ràng buộc chứ không xác nhận, nên câu hỏi xây dựng bậc hai có đúng
hay không phải được trả lời bằng phép đo. Bài trả lời bằng ba thí nghiệm trên
ba họ mô hình, và mục này đọc chúng theo một trục duy nhất: mỗi phép đo lấy cái
gì làm nền để nói kết quả là đúng.

Thí nghiệm thứ nhất chạy trên mô hình ngôn ngữ. Bài đo tương tác giữa các token
của \tn{lời nhắc}{prompt} trong tám mô hình Gemma 3, gồm bốn cỡ 1B, 4B, 12B và 27B tham số,
mỗi cỡ hai bản, bản \tn{tiền huấn luyện}{pretraining} và bản đã tinh chỉnh theo chỉ dẫn. Bộ lời
nhắc có 21 câu, in đầy đủ trong phần phụ lục của bài, mỗi câu sinh ra câu trả
lời dài tối đa 100 token~\autocite{p23metagame}. Đại lượng được báo cáo là
Recall@K, và định nghĩa của nó là chỗ mục này quan tâm: nó đo xem những cặp
token do người gán nhãn xếp hạng thế nào so với những tương tác mà mô hình cho
là quan trọng nhất, với $K$ chạy từ 0 tới 15 phần trăm số cặp có hướng. Người
gán nhãn ở đây là chính các tác giả: trong danh sách 21 lời nhắc, những cặp
token mà họ coi là có tương tác được in đậm, mỗi lời nhắc từ hai tới bốn cặp.

Vậy nền của thí nghiệm thứ nhất là một phán đoán của con người về những cặp
token lẽ ra phải tương tác với nhau. Theo cách chia mà
mục~\ref{sec:ch01-do-trung-thuc-va-tinh-hop-ly} đặt ra, đó là một phép đo ở
phía \tn{tính hợp lý}{plausibility}, không phải ở phía \tn{độ trung thực}{faithfulness}: nó hỏi
lời giải thích có khớp với điều người đọc chờ đợi hay không, chứ không hỏi nó
có khớp với điều mô hình đã làm hay không. Bài không nói ngược lại; nó không
gọi thí nghiệm này là một phép đo độ trung thực ở bất cứ chỗ nào.

Thí nghiệm thứ hai chạy trên năm bộ mã hóa ảnh và văn bản: một bản CLIP, hai
bản SigLIP-2 và hai bản MetaCLIP-2, trong đó bản MetaCLIP-2 thứ hai chỉ dùng để
vẽ minh họa ở độ phân giải cao chứ không góp số vào bảng. Đại lượng được báo cáo là một chỉ số nhận diện tương
tác mượn từ một bài báo khác, chạy trên tập kiểm định của ImageNet-1k: bài lấy
mười nhãn lớp, mỗi nhãn 50 ảnh, ghép thành mười lượt thử bốn nhãn, rồi thêm dần
từng nhãn một để có bốn tình huống mỗi lượt~\autocite{p23metagame}. Phép đo
hỏi tương tác có rơi đúng vào vật mà nhãn ghi hay không.

Kết quả của thí nghiệm thứ hai có một dạng đáng ghi lại. Với CLIP, chỉ số của
phương pháp dựa trên \tn{cơ chế chú ý}{attention mechanism} là 0.77 khi ảnh có một vật và tụt xuống 0.25
khi có bốn vật, còn bản Meta của chính nó thì đi ngược, từ 0.84 lên 0.90. Với
MetaCLIP-2 cũng vậy: 0.84 tụt xuống 0.25, còn bản Meta giữ ở 0.85 và
0.86~\autocite{p23metagame}. Trong bảng ấy có năm dòng phương pháp bậc nhất,
thuộc ba phương pháp khác nhau chạy trên hai mô hình, và cả năm dòng đều dừng
đúng ở 0.25 tại cột bốn vật. Con số 0.25 gợi mức bốc
thăm, nhưng chỉ số này được định nghĩa trong một bài báo mà sách không đọc, và
bài đang đọc không nói mức bốc thăm của nó là bao nhiêu, nên sách ghi lại chỗ
trùng ấy mà không phát biểu nó nghĩa là gì. Bài cũng ghi rằng một tổ hợp mô
hình và phương pháp đã bị loại khỏi bảng vì cho đầu ra không đọc được, nên bảng
không phải là toàn bộ những gì đã chạy.

Thí nghiệm thứ ba chạy trên một mô hình sinh ảnh từ văn bản và đo bằng một
nhiệm vụ \tn{phân đoạn}{segmentation}. Trên Pascal VOC, tập kiểm định gồm 1\,449 ảnh, trong đó 927
ảnh có một khái niệm được chú thích và 522 ảnh có từ hai tới năm; trên MS COCO
là 5\,000 ảnh chú thích toàn cảnh, còn 4\,660 ảnh đủ điều kiện sau khi lọc; và
trên ImageNet-Seg là 4\,276 mặt nạ nhị phân, giữ lại 3\,535~\autocite{p23metagame}.
Ở nhiệm vụ một khái niệm trên Pascal VOC, độ chính xác tăng từ 0.832 lên 0.903
khi thêm xây dựng bậc hai vào phương pháp gốc.

Ba cái nền ấy đặt cạnh nhau thì thấy chúng cùng một loại. Cái thứ nhất là phán
đoán của người, cái thứ hai là nhãn lớp của bộ dữ liệu, cái thứ ba là mặt nạ
của bộ dữ liệu. Cả ba đều là câu trả lời cho câu hỏi trong đầu vào có cái gì và
nằm ở đâu, và không cái nào là câu trả lời cho câu hỏi mô hình đã dùng cái gì.
Đó là đúng khoảng cách mà định nghĩa~\ref{def:ch01-do-trung-thuc} tách ra, và
là lý do vì sao ba bảng kết quả ấy, dù đọc kỹ đến đâu, cũng không cho biết
meta-attribution có trung thực hay không.

Một con số cuối cùng của phần thí nghiệm đáng in ra vì nó đo cái giá. Bài ghi
rằng toàn bộ thí nghiệm chạy trên một cụm có tám GPU H100 trong khoảng 15 ngày,
và các tác giả ước rằng những lần chạy thử và những lần thất bại tiêu tốn thêm
chừng ấy nữa~\autocite{p23metagame}. Chỗ đứng của con số ấy là ở
mục~\ref{sec:ch11-khong-co-tang-day}, khi câu hỏi trở thành bước tiếp theo tốn
thêm bao nhiêu.
\index{đánh giá!nền của phép đo|)}%
